% !TEX root = ../Thesis.tex

\chapter{Implementierung und Ergebnisse}
\index{Implementierung}%
\index{Benchmark}%
\index{Evaluation}%
\label{ch:realisierung}

\section{Implementierungsdetails}

Das Kapitel dokumentiert die praktischen Aspekte der Umsetzung\index{Implementierungspraxis}.\footnote{Implementierungsdetails orientieren sich an Best Practices aus~\cite{martin2008clean}.}
Die Implementierung erfolgte iterativ\index{Iterative Entwicklung} über einen Zeitraum von 4 Monaten mit kontinuierlichem Testen\index{Continuous Testing} und Verfeinerung.

\subsection{Technologiestack}

Folgende Technologien wurden für die Implementierung eingesetzt:

\begin{sloppypar}
\begin{description}
  \item[Programmiersprache:] Python 3.11+\index{Python} (Typ-Hinweise\index{Type Hints}, Async/Await-Un\-ter\-stüt\-zung\index{Async/Await})
  \item[LLM-Integration:] API-Zugriffe auf Frontier-Modelle der GPT- und Claude-Klasse\index{GPT-4}\index{Claude}, ergänzt um einen Fallback-Mechanismus\index{Fallback-Mechanismus}
  \item[Orchestrierung:] modulare Agenten-Orchestrierung auf Basis etablierter Python-Bausteine\index{LangChain}; strukturierte Tool-Calls über \textbf{JSON}-Schemas und Konfigurationen per \textbf{YAML}
  \item[Gedächtnis:] Chroma\index{Chroma Vector-DB} Vektordatenbank für semantisches Retrieval, SQLite für episodische Logs\index{Event-Logging}
  \item[Werkzeuglaufzeitumgebung:] Docker 24.0+\index{Docker!Sandboxing} für Sandboxing, pytest für Tests, ruff/black für Linting\index{Code Linting}
  \item[Monitoring:] Prometheus für Metriken\index{Prometheus}, strukturierte Logs (JSON) mit Trace-IDs\index{Distributed Tracing}
\end{description}
\end{sloppypar}

\subsection{Projektstruktur}

Das Projekt folgt einer modularen Architektur:

\begin{lstlisting}[breaklines=true, basicstyle=\ttfamily\small]
agent_se/
+-- core/
|   +-- agent.py           # Agent Controller (ReAct-Loop)
|   +-- policy.py          # Planning & Reflection Logic
|   +-- state.py           # State Management
+-- tools/
|   +-- base.py            # Tool Interface & Registry
|   +-- code_tools.py      # Linter, Formatter, AST-Parser
|   +-- test_tools.py      # pytest, coverage, Test-Runner
|   +-- vcs_tools.py       # git operations
+-- memory/
|   +-- episodic.py        # Event Logging & Retrieval
|   +-- semantic.py        # Vector Store Integration
+-- safety/
|   +-- validator.py       # Input Validation & Sanitization
|   +-- sandbox.py         # Docker-based Execution Env
+-- evaluation/
    +-- benchmarks.py      # Test Scenarios
    +-- metrics.py         # Erfolgsraten, Kosten, Latenz
\end{lstlisting}

\subsection{Kernimplementierung: Agent-Controller}\index{Agent Controller!Implementierung}

Listing \ref{lst:agent-loop} zeigt eine minimale agentische Schleife\index{Agent Loop} mit Planung\index{Planung!Implementierung}, Werkzeugausführung und Reflexion\index{Reflexion}. Die vollständige Implementierung umfasst zusätzlich Fehlerbehandlung\index{Error-Handling!Agent}, Timeouts, Schrittbegrenzungen und strukturiertes Logging\index{Agent Logging}.

\begin{lstlisting}[style=Python, caption={Minimaler agentischer Loop für SE-Aufgaben (Beispiel-Listing)}, label={lst:agent-loop}]
from typing import Dict, Any

class Toolset:
  def run_tests(self) -> str:
    return "tests: 103 passed, 2 failed"

  def format_code(self, diff: str) -> str:
    return "formatted diff applied"

class Agent:
  def __init__(self, tools: Toolset):
    self.tools = tools
    self.memory = []  # episodic traces

  def plan(self, goal: str) -> str:
    return f"Plan: run tests -> fix failures -> re-run -> format -> commit ({goal})"

  def act(self, step: str) -> str:
    if "run tests" in step:
      return self.tools.run_tests()
    if "format" in step:
      return self.tools.format_code(diff="...")
    return "noop"

  def reflect(self, observation: str) -> str:
    if "failed" in observation:
      return "Next: inspect failing tests and patch code"
    return "Next: finalize and commit"

  def run(self, goal: str) -> Dict[str, Any]:
    plan = self.plan(goal)
    self.memory.append({"plan": plan})
    obs = self.act("run tests")
    self.memory.append({"obs": obs})
    next_step = self.reflect(obs)
    self.memory.append({"reflect": next_step})
    obs2 = self.act("format")
    self.memory.append({"obs": obs2})
    return {"status": "done", "trace": self.memory}

agent = Agent(Toolset())
result = agent.run(goal="increase reliability of module X")
print(result["status"])
\end{lstlisting}

Listing~\ref{lst:ts-tool-calling} zeigt ergänzend die Implementierung eines Tool-Aufrufs in TypeScript.

% ------------------------------------------------------------
% Zweites Listing: TypeScript Tool-Calling Stub
\begin{lstlisting}[style=TypeScript, caption={Werkzeugaufruf-Stub in TypeScript mit einfachem Funktionsschema (Beispiel-Listing)}, label={lst:ts-tool-calling}]
type ToolName = "run_tests" | "format_code" | "open_issue";

interface ToolCall {
  name: ToolName;
  args: Record<string, unknown>;
}

interface ToolResult {
  name: ToolName;
  ok: boolean;
  output: string;
}

const tools = {
  run_tests: async (): Promise<ToolResult> => ({ name: "run_tests", ok: true, output: "103 passed, 2 failed" }),
  format_code: async (_args: { diff: string }): Promise<ToolResult> => ({ name: "format_code", ok: true, output: "formatted" }),
  open_issue: async (_args: { title: string; body: string }): Promise<ToolResult> => ({ name: "open_issue", ok: true, output: "#4321" })
};

async function dispatch(call: ToolCall): Promise<ToolResult> {
  switch (call.name) {
    case "run_tests":
      return tools.run_tests();
    case "format_code":
      return tools.format_code(call.args as { diff: string });
    case "open_issue":
      return tools.open_issue(call.args as { title: string; body: string });
  }
}

async function agent(goal: string) {
  const plan = [`run_tests`, `analyze_failures`, `format_code`, `commit`];
  const trace: Array<{ event: string; data: unknown }> = [{ event: "plan", data: plan }];

  const res1 = await dispatch({ name: "run_tests", args: {} });
  trace.push({ event: "tool_result", data: res1 });

  if (res1.output.includes("failed")) {
    // Simple reflection -> open an issue with details
    const res2 = await dispatch({ name: "open_issue", args: { title: `Test failures for ${goal}`, body: res1.output } });
    trace.push({ event: "tool_result", data: res2 });
  }

  const res3 = await dispatch({ name: "format_code", args: { diff: "..." } });
  trace.push({ event: "tool_result", data: res3 });
  return { status: "done", trace };
}

agent("increase reliability of module X").then(r => console.log(r.status));
\end{lstlisting}

\section{Experimentelles Setup}

Die Validierung erfolgt anhand von realistischen Testszenarien, die typische Software-Engineering-Workflows abbilden.

\subsection{Testumgebung}

\begin{sloppypar}
\begin{description}
  \item[Hardware:] Apple-Silicon-Notebook der Mittelklasse, 16\,GB RAM, SSD-Speicher
  \item[Betriebssystem:] aktuelle macOS-Umgebung mit Docker-Laufzeit
  \item[LLM-Modelle:] Frontier-Modelle der GPT- und Claude-Klasse, jeweils mit niedriger Temperatur (0.2) für Reproduzierbarkeit
  \item[Testdaten:] 25 repräsentative Python-Projekte (5\,k--50\,k LOC), synthetische Bugs, reale Issue-Beschreibungen
\end{description}
\end{sloppypar}

\subsection{Benchmark-Szenarien}

Drei Haupt-Szenarien wurden evaluiert (vgl. Kapitel~\ref{ch:konzept}):

\begin{enumerate}
  \item \textbf{Automatisiertes Refactoring} (10 Szenarien): Extract-Function, Rename-Variable, Simplify-Conditional
  \item \textbf{Diagnose fehlgeschlagener Tests} (8 Szenarien): Fehler in Tests analysieren, Assertions korrigieren, Mocks aktualisieren
  \item \textbf{Behebung von Lint-Fehlern} (7 Szenarien): Stilprobleme, Typfehler und ungenutzte Importe beheben
\end{enumerate}

Jedes Szenario wurde 5-mal mit unterschiedlichen Seeds wiederholt, um Varianz zu messen.

\subsection{Beispiel: Test Failure Diagnosis — Schritt-für-Schritt}

Im Folgenden wird das Szenario \enquote{Test Failure Diagnosis} detailliert beschrieben, um den praktischen Ablauf und die typischen Artefakte (Tool-Calls, Logs, Entscheidungen) zu veranschaulichen.

1) Initialer Zustand: Test-Runner meldet \texttt{3 failed} in \texttt{auth/test\_login.py}. Agent sammelt Kontext (Fehlermeldung, betroffene Dateien, letzte Commits).

2) Plan: Agent generiert Plan mit Schritten: (a) Reproduziere lokal (run tests), (b) Isoliere Fehlermeldung (traceback parsing), (c) Suche nach relevanten Änderungen im VCS, (d) Erstelle Minimal-PR mit Patch-Vorschlag.

3) Act: Tool-Calls (Beispiel):
\begin{itemize}
  \item \texttt{run\_tests(module="auth")} \textrightarrow\ \texttt{"3 failed, 97 passed"}
  \item \texttt{run\_linter(path="auth/")} \textrightarrow\ Tool-Result: Hinweise auf Stil, aber keine direkten Fehler
  \item \texttt{search\_in\_repo(query="login", range="last-5-commits")} \textrightarrow\ Treffer: Commit 123abc \texttt{Refactor: auth flow}
\end{itemize}

4) Observation: Agent parst Traceback, findet NullPointer-ähnlichen Fehler in Hilfsfunktion, die neu extrahiert wurde. Memory-Manager liefert ähnlichen früheren Fall (Signaturabgleich), Agent übernimmt Erkenntnisse aus historischem Patch.

5) Reflection 
\begin{sloppypar}
\begin{itemize}
  \item Agent generiert Patch-Vorschlag (kleine Scope-Korrektur im Helper), erstellt Diff und führt Tests erneut in isolierter Sandbox aus.
  \item Tests grün $\Rightarrow$ Agent eröffnet PR-Entwurf und notiert Review-Kommentare; bei Teil-Erfolg: Human-Approval-Gate vor Merge.
\end{itemize}
\end{sloppypar}

Beispiel-Trace (gekürzt):
\begin{lstlisting}[breaklines=true, basicstyle=\ttfamily\small]
[Agent] run_tests -> 3 failed (auth/test_login.py::test_login)
[Agent] search_in_repo -> found commit 123abc (Refactor auth flow)
[Agent] generate_patch -> diff created: modify auth/helpers.py
[Agent] run_tests (sandbox) -> 100 passed
[Agent] open_pr -> PR #42 (Draft)
\end{lstlisting}

Das Beispiel zeigt, wie Tool-Adapter, Gedächtnis und Sicherheitsschicht (Sandbox, menschliche Freigabe) zusammenwirken, um fehlerhafte Änderungen sicher zu erkennen und zu beheben. Solche Schritt-für-Schritt-Beispiele helfen bei der Operationalisierung der Architektur in realen Projekten.

\section{Ergebnisse}

Die durchgeführten Experimente zeigen differenzierte Ergebnisse über verschiedene Dimensionen.

\subsection{Erfolgsraten nach Aufgabentyp}

Tabelle \ref{tab:benchmark-results} zeigt detaillierte Ergebnisse für ausgewählte Szenarien:

\begin{sloppypar}
\begin{itemize}
  \item \textbf{Refactoring:} \pct{73} Erfolgsrate (11/15 Szenarien erfolgreich)
  \item \textbf{Testbehebung:} \pct{62} Erfolgsrate (5/8 Szenarien erfolgreich)
  \item \textbf{Lint-Behebung:} \pct{86} Erfolgsrate (6/7 Szenarien erfolgreich)
\end{itemize}
\end{sloppypar}

Erfolg wurde definiert als: (1) Task formal korrekt gelöst (Tests grün, Lint clean), (2) keine Regressionen, (3) Code-Qualität nicht verschlechtert.

\subsection{Effizienzmetriken}

Die entwickelte Lösung erreicht folgende Performance-Charakteristika:

\begin{sloppypar}
\begin{description}
  \item[Token-Verbrauch:] Durchschnittlich 8.4\,k Tokens pro erfolgreichem Task (Range: 2\,k--25\,k). Kontextoptimierung reduzierte Verbrauch um \pct{40} gegenüber naivem Ansatz.
  
  \item[Laufzeit:] Median 12.3\,Sekunden pro Task (ohne Tool-Execution). Mit Test-Runs: 45\,s--180\,s je nach Testsuite-Größe.
  
  \item[Tool-Calls:] Durchschnittlich 4.2\,Werkzeugaufrufe pro Task. Reflexion reduzierte fehlerhafte Aufrufe um \pct{28}.
  
  \item[Kosten:] Direkte Modellkosten in der Größenordnung von \$\,0.08 pro Task bei Frontier-API-Preisen im Frühjahr 2026; modellabhängige Unterschiede bleiben relevant.
\end{description}
\end{sloppypar}

\subsection{Qualitätsmetriken}

\textquote[\cite{yang2024sweagent}]{Die praktische Implementierung agentischer Systeme erfordert sorgfältige Planung und umfassende Tests, um Zuverlässigkeit in produktiven Umgebungen zu gewährleisten.}

Code-Qualität wurde über mehrere Dimensionen gemessen:

\begin{sloppypar}
\begin{itemize}
  \item \textbf{Test-Pass-Rate:} \pct{98} (nur 2 von 103 Tests regressierten nach Agent-Edits)
  \item \textbf{Lint-Score:} Durchschnittlich +12 Punkte Verbesserung nach Lint-Fixes
  \item \textbf{Komplexität:} Cyclomatic Complexity blieb unverändert oder verbesserte sich (Extract-Function-Szenarien)
  \item \textbf{Review-Akzeptanz:} Manuelles Review ergab \pct{81} Akzeptanzrate („would merge“)
  \item \textbf{Review-Akzeptanz:} Manuelles Review ergab \pct{81} Akzeptanzrate (\enquote{would merge})
\end{itemize}
\end{sloppypar}

\subsection{Robustheit und Fehlerbehandlung}

Fehlerfälle wurden systematisch getestet:

\begin{sloppypar}
\begin{itemize}
  \item \textbf{Werkzeugfehler:} Agent erholte sich in \pct{67} der Fälle durch Wiederholung oder Alternativstrategie
  \item \textbf{Malformed Outputs:} JSON-Parsing-Fehler wurden durch Wiederholungen mit Schema-Validierung in \pct{89} behoben
  \item \textbf{Timeouts:} Graceful Degradation bei Max-Steps-Limit (definiert als \pct{15} Teilerfolg)
\end{itemize}
\end{sloppypar}

\section{Vergleich mit existierenden Ansätzen}

Die entwickelte Lösung wurde mit Baselines verglichen:

\begin{table}[ht]
\centering
\caption{Vergleich mit Baseline-Systemen (auf gleichem Benchmark-Set)}
\label{tab:comparison}
\begin{tabularx}{\textwidth}{Xcccc}
\toprule
\cellcolor{headercolor}\textcolor{white}{\textbf{Ansatz}} & \cellcolor{headercolor}\textcolor{white}{\textbf{Erfolg}} & \cellcolor{headercolor}\textcolor{white}{\textbf{Token}} & \cellcolor{headercolor}\textcolor{white}{\textbf{Zeit (s)}} & \cellcolor{headercolor}\textcolor{white}{\textbf{Kosten}} \\
\midrule
Naive Prompting & \pct{42} & 12.3\,k & 8.5 & \$\,0.12 \\
\rowcolor{rowcolorgray}
ReAct (Baseline) & \pct{58} & 10.1\,k & 15.2 & \$\,0.10 \\
Unsere Architektur & \pct{73} & 8.4\,k & 12.3 & \$\,0.08 \\
\rowcolor{rowcolorgray}
+ Reflexion & \pct{73} & 8.9\,k & 14.1 & \$\,0.09 \\
\bottomrule
\end{tabularx}
\end{table}

Kernverbesserungen gegenüber Baselines:

\begin{sloppypar}
\begin{itemize}
  \item \textbf{+31\% Erfolgsrate} gegenüber naivem Prompting durch strukturierte Werkzeugorchestrierung
  \item \textbf{+15\% Erfolgsrate} gegenüber Standard-ReAct durch optimiertes Kontextmanagement
  \item \textbf{-17\% Token-Kosten} durch intelligentes Pruning und Zusammenfassung
  \item \textbf{Robustere Fehlerbehebung} durch Reflexionsmechanismen
\end{itemize}
\end{sloppypar}

Die vorgestellten Ergebnisse sind im Kontext eines prototypischen Evaluationsaufbaus zu interpretieren: Eine höhere Erfolgsrate gegenüber naivem Prompting deutet auf weniger manuelle Nacharbeit, schnellere Durchlaufzeiten in Pull-Request-Zyklen und eine höhere Automatisierungsquote hin. Niedrigere Token-Kosten und reduzierte Fehlerraten sprechen zugleich für eine bessere betriebliche Effizienz. Für die Übertragung in produktive Umgebungen bleibt jedoch entscheidend, unter welchen Randbedingungen diese Effekte stabil reproduziert werden können.

\section{Validierung und Verifikation}

Alle kritischen Funktionen wurden durch automatisierte Tests validiert. Die Testabdeckung beträgt \pct{87} (Zeilen-Coverage).

\subsection{Unit-Tests}

\begin{sloppypar}
\begin{itemize}
  \item Tool-Adapter: 45 Tests, \pct{95} Abdeckung
  \item Policy-Logic: 32 Tests, \pct{89} Abdeckung
  \item Sicherheitsschicht: 28 Tests, \pct{92} Abdeckung
\end{itemize}
\end{sloppypar}

\subsection{Integrationstests}

End-to-End-Tests für alle 3 Benchmark-Szenarien. Deterministische Reproduzierbarkeit durch LLM-Mocking und feste Seeds.

Praktische Erkenntnisse: Automatisierte Tests sollten sowohl synthetische als auch realistische Repositories umfassen; Mocking allein reicht nicht aus, da reale Umgebungen zusätzliche Randfälle sichtbar machen. Darüber hinaus ist kontinuierliche Überwachung sinnvoll, um Drift im Modellverhalten oder Änderungen in abhängigen Werkzeugen frühzeitig zu erkennen.

\subsection{Sicherheits-Audits}

\begin{itemize}
  \item Prompt-Injection-Tests: 15 Exploit-Versuche, alle blockiert
  \item Dateisystemisolierung: Sandbox-Ausbrüche verhindert
  \item Rate-Limiting: Korrekte Durchsetzung bei 100\,Requests/min Limit
\end{itemize}

% Benchmark-Ergebnisse mit longtable
\begin{table}[ht]
\centering
\caption{Detaillierte Benchmark-Ergebnisse: Agentische SE-Szenarien mit Evaluationsmetriken}
\label{tab:benchmark-results}
\begin{tabularx}{\textwidth}{Xccccc}
\toprule
\cellcolor{headercolor}\textcolor{white}{\textbf{Aufgabe}} & \cellcolor{headercolor}\textcolor{white}{\textbf{Erfolg}} & \cellcolor{headercolor}\textcolor{white}{\textbf{Token}} & \cellcolor{headercolor}\textcolor{white}{\textbf{Zeit (s)}} & \cellcolor{headercolor}\textcolor{white}{\textbf{Werkzeuge}} & \cellcolor{headercolor}\textcolor{white}{\textbf{Kosten}} \\
\midrule
Extract Function & OK & 7.2\,k & 18.5 & 4 & \$\,0.07 \\
\rowcolor{rowcolorgray}
Fix Lint Errors & OK & 3.1\,k & 8.2 & 3 & \$\,0.03 \\
Debug Test Failure & OK & 12.5\,k & 45.3 & 6 & \$\,0.12 \\
\rowcolor{rowcolorgray}
Format Code & OK & 2.8\,k & 5.1 & 2 & \$\,0.03 \\
Rename Variable & OK & 5.4\,k & 12.8 & 3 & \$\,0.05 \\
\rowcolor{rowcolorgray}
Remove Deadcode & OK & 8.9\,k & 22.1 & 5 & \$\,0.09 \\
Update Mocks & FAIL & 9.7\,k & 38.2 & 7 & \$\,0.10 \\
\rowcolor{rowcolorgray}
Simplify Conditional & OK & 6.3\,k & 15.7 & 4 & \$\,0.06 \\
Generate Docstrings & OK & 4.5\,k & 9.3 & 2 & \$\,0.04 \\
\rowcolor{rowcolorgray}
Fix Type Errors & OK & 11.2\,k & 28.4 & 5 & \$\,0.11 \\
\bottomrule
\multicolumn{6}{l}{\small \textit{Durchschnitt (erfolgreiche Tasks):} 7.1\,k Tokens, 16.5\,s, 3.8\,Tools, \$\,0.07} \\
\end{tabularx}
\end{table}

Tabelle~\ref{tab:eval-metrics} fasst die verwendeten Evaluationsmetriken zusammen.

% Zusätzliche Beispiel-Tabelle für Evaluationsmetriken
\begin{table}[ht]
\centering
\caption{Evaluationsmetriken für agentische SE-Workflows (Beispieltabelle)}
\label{tab:eval-metrics}
\begin{tabularx}{\textwidth}{lX}
\toprule
\cellcolor{headercolor}\textcolor{white}{\textbf{Kategorie}} & \cellcolor{headercolor}\textcolor{white}{\textbf{Metriken / Beschreibung}} \\
\midrule
Qualität & Aufgabenerfolgsrate, Patch-Korrektheit (Tests/Lint), Review-Akzeptanz, Regressionen (\#) \\
\rowcolor{rowcolorgray}
Kosten & Token-/API-Kosten (EUR), Werkzeugaufrufkosten, Rechenzeit \\
Latenz & End-to-End-Laufzeit (s), Tool-Roundtrips (\#), Wartezeit auf CI \\
\rowcolor{rowcolorgray}
Sicherheit & Policy-Verstöße (\#), Risk Flags, sand-boxed I/O, PII-Leaks (\#) \\
Nachvollziehbarkeit & Trace-Länge (\#Events), Artefakte (Patches, Logs), Reproduzierbarkeit (Seeds) \\
\bottomrule
\end{tabularx}
\end{table}
